\section{Paths through the principal stratum}\label{sec:paths}

Before comparing compact spaces, we identify the central metric that
will be their limit. Density of the stable locus is not sufficient:
a dense Riemannian stratum can have a larger intrinsic distance.
We prove equality of distances by perturbing paths and controlling
the short joins between their pieces.

\subsection{Short paths in a fixed semialgebraic set}
We use the following elementary consequence of semialgebraic
triangulation and uniform monotonicity
\cite[pp.~116,130,133]{vandendries}.

\begin{lemma}\label{lem:bounded-paths}
A nonempty bounded connected semialgebraic set \(S\subset\R^N\),
not necessarily closed, has finite intrinsic diameter for
Euclidean path length. More precisely, every two points can be
joined in \(S\) by a continuous semialgebraic path of length at
most a constant \(C_S\).
\end{lemma}

\begin{proof}
Triangulate the compact closure compatibly with \(S\).
Its inverse image is a finite union of relative open simplices.
Choose their barycenters and connect two when one included simplex
is a face of the other. The connecting segment lies in the larger
open simplex except at its endpoint and stays in the union.
The resulting finite incidence graph is connected. Otherwise one
component and its complement would be relatively open and closed:
a relative limit between included cells can only occur along a
face incidence, which would connect the components.

Choose a spanning tree. Join each variable endpoint to its cell
barycenter and then follow the tree between barycenters. Transport
the resulting polygonal paths through the semialgebraic
triangulation map. There are finitely many endpoint-cell choices,
so this is a semialgebraic family of continuous paths. Cell decomposition and uniform finiteness, applied with the two
endpoints as parameters, bound the number of monotonicity intervals
of each coordinate by a common integer. Every coordinate stays in the
bounded range of \(S\). The sum of its coordinate variations
therefore gives a uniform finite bound for path length.
\end{proof}

For a positive-dilation-invariant set, applying the lemma in its
unit ball gives joining paths of length \(C_S\rho\) in its
radius-\(\rho\) ball. This is a statement about a fixed set; it
does not assert a uniform constant for arbitrary degenerating
families.

\subsection{The simple zero quotient}
Let \(W\) be a full doubled-quiver representation, with compact
group \(K\) a connected product of unitary groups modulo an
ineffective common scalar. Put
\[
 A_0=\{\mu_{\C}=0,\ \mu_{\R}=0\},\qquad Q_0=A_0/K.
\]
Assume that the complex zero fibre contains a module-simple
representation. Write \(A_0^s\) for the balanced simple locus and
\(Q_0^s=A_0^s/K\).
We use the finite intrinsic theorem
\cite[Section~5]{wuebben14}: \(Q_0\), with the infimum of
finite-stratum path lengths, is a proper geodesic space whose
metric topology is the quotient topology. We also use its full
normal metric slice and normal-simple results
\cite[Sections~4,6]{wuebben14}. These retain the complete normal
representation and require the simple-zero hypothesis.

\begin{lemma}\label{lem:simple-paths}
There is a constant \(C_W\) such that every two points of
\(A_0^s\) of norm at most \(\rho\) can be joined in \(A_0^s\),
within that ball, by a path of length at most \(C_W\rho\).
At any point of \(Q_0\), nearby simple points admit simple joining
paths whose length tends to zero and which stay in any fixed
sufficiently small neighborhood.
\end{lemma}

\begin{proof}
Crawley--Boevey's theorem \cite[Theorem~1.2]{crawleyboevey},
under the existence of a simple representation, makes the complex
zero fibre irreducible and generically simple. Its simple quotient
is a connected smooth open complex variety. Kempf--Ness identifies
this with \(Q_0^s\). The compact-group quotient
\(A_0^s\to Q_0^s\) is proper with connected fibres. A separation
of its source would induce a separation of its target, by the
closed-map property and connected fibres. Hence \(A_0^s\) is connected.

Simplicity is semialgebraic: for each proper subdimension vector,
the existence of invariant subspaces is a polynomial incidence
condition on a product of Grassmannians, followed by projection.
There are finitely many such vectors. The balanced moment equations
are real polynomial, and \(A_0^s\) is invariant under positive
dilation. The map \(x\mapsto x/\max(1,\|x\|)\) retracts it onto its
unit-ball part. That part is connected and bounded, so
Lemma~\ref{lem:bounded-paths} and scaling prove the first assertion.
The zero-dimensional case uses constant paths.

For the second assertion take the full normal slice at a balanced
representative of the chosen point. Its stabilizer is connected,
and its full normal representation again has a module-simple point
by the normal-simple theorem. Representatives of nearby points
tend, after unitary adjustment, to zero in this slice.
A closed zero-level representation has scalar stabilizer exactly
when it is simple; the slice preserves this condition.
Apply the first assertion in the full normal representation.
The bounded metric comparison of the slice takes its \(C\rho\)
joins to joins of length at most \(C'\rho\) in the ambient quotient.
Their images lie in any prescribed fixed neighborhood for small
\(\rho\).
\end{proof}

Projection cannot increase the lengths of the balanced paths.
Semialgebraic paths are piecewise smooth after finite subdivision,
and within the simple stratum they may be smoothed with arbitrarily
small additional length.

\begin{proposition}\label{prop:principal-completion}
The intrinsic Riemannian distance of \(Q_0^s\) equals the restriction
of the full finite-stratum distance of \(Q_0\).
Its metric completion is \(Q_0\).
\end{proposition}

\begin{proof}
One inequality follows because every simple path is admissible
in the full quotient. For the reverse, take simple endpoints and
a finite-stratum path of length within \(\epsilon\) of their full
distance. Truncate an arbitrarily small total length at its finitely
many stratum junctions. Cover the remaining compact smooth pieces
by normal slice charts with relative squared-metric errors less
than a chosen \(\delta>0\), and subdivide.

In a full slice the given stratum is \(T\times\{0\}\), where the
group acts trivially on \(T\). Choose balanced simple normal points
\((u_j,z_j)\to0\). Translation in \(T\) removes \(u_j\), without
changing the moment map or stabilizer. Replace a piece
\((u(t),0)\) by \((u(t),z_j)\). Its product-metric speed is exactly
the old speed. Comparing both curves with that same product metric
increases length by at most
\(\sqrt{(1+\delta)/(1-\delta)}\). Their endpoints converge.

Use Lemma~\ref{lem:simple-paths} to join the perturbed endpoints
at each subdivision point. At an original singular junction,
first choose a neighborhood where the joining cost is below an
assigned tolerance, then truncate far enough toward the junction
to place both endpoints there. Small simple perturbations remain
there. The same construction joins to the original simple endpoints.
There are finitely many joins; allocate their total cost below
\(\epsilon\), and choose the chart errors so their total length
inflation is below \(\epsilon\).
The resulting simple path has length at most the full distance
plus \(3\epsilon\). Let \(\epsilon\downarrow0\).

The simple locus is dense in the quotient topology, hence in the
finite metric topology. Properness gives completeness of \(Q_0\),
so it is the asserted completion.
\end{proof}

This argument explains why lower-dimensional strata create no
new shortcuts between stable points: every such path can be
approximated with its length, not merely with its endpoints.
